%───────────────────────────────────────────────────────────────────────────────
% TABLES.TEX — Estilo global de tablas con colores semánticos Catppuccin
%───────────────────────────────────────────────────────────────────────────────
% Consume aliases del dominio table-* definidos en catppuccin-palette.tex.
% No usa Ctp* directamente ni define colorlets propios.
%
% ENTORNOS DISPONIBLES:
%   acadtable   → tabla flotante (fuera de tcolorbox)
%   boxtab      → tabla no-flotante para DENTRO de tcolorbox (en tcolorbox.tex)
%   fittable    → reescala automática si la tabla excede el ancho disponible
%───────────────────────────────────────────────────────────────────────────────


% ─── Espaciado global de booktabs ────────────────────────────────────────────
\setlength{\heavyrulewidth}{0.08em}
\setlength{\lightrulewidth}{0.04em}
\setlength{\cmidrulewidth}{0.03em}
\renewcommand{\arraystretch}{1.25}

% ─── Alineación vertical: m (middle) en lugar de p (top) ────────────────────
% tabularx usa p{} por defecto → desalineación vertical al mezclar con l/c/r.
\renewcommand{\tabularxcolumn}[1]{m{#1}}

% ─── Tipos de columna ────────────────────────────────────────────────────────
% L/C/R: ancho fijo con wrapping y centrado vertical
\newcolumntype{L}[1]{>{\raggedright\arraybackslash}m{#1}}
\newcolumntype{C}[1]{>{\centering\arraybackslash}m{#1}}
\newcolumntype{R}[1]{>{\raggedleft\arraybackslash}m{#1}}

% Y/Z: columnas proporcionales para tabularx (centrada / derecha)
% X ya existe en tabularx (izquierda, proporcional, ahora con m{})
\newcolumntype{Y}{>{\centering\arraybackslash}X}
\newcolumntype{Z}{>{\raggedleft\arraybackslash}X}

% ─── Color de reglas booktabs ─────────────────────────────────────────────────
\AtBeginDocument{\arrayrulecolor{table-rule}}

% ─── Encabezado de columna ───────────────────────────────────────────────────
\newcommand{\thead}[1]{%
  \bfseries\color{table-head-text}#1%
}

% Colorear la fila de encabezado completa vía \rowcolor
\newcommand{\headerrow}{\rowcolor{table-head}}


% ─── fittable: reescala si excede el ancho disponible ────────────────────────
% Funciona tanto en texto normal (\textwidth) como dentro de tcolorbox
% (\linewidth). Usa \linewidth en ambos contextos: fuera de una caja
% \linewidth == \textwidth; dentro, \linewidth es el ancho de la caja.
%
% Uso:
%   \begin{fittable}
%     \begin{tabular}{llll} ... \end{tabular}
%   \end{fittable}
\newsavebox{\fittablebox}
\NewDocumentEnvironment{fittable}{}{%
  \begin{lrbox}{\fittablebox}%
}{%
  \end{lrbox}%
  \ifdim\wd\fittablebox>\linewidth
    \resizebox{\linewidth}{!}{\usebox{\fittablebox}}%
  \else
    \usebox{\fittablebox}%
  \fi
}


% ─── acadtable: tabla flotante (SOLO fuera de tcolorbox) ─────────────────────
% Produce \begin{table}[placement] → flotante.
% DENTRO de un entorno tcolorbox (definicion, ejemplo, etc.) usar boxtab
% de config/tcolorbox.tex en su lugar.
%
% Uso:
%   \begin{acadtable}[htbp]{tab:label}{Caption text}
%     \begin{tabularx}{\textwidth}{l Y Z}
%       ...
%     \end{tabularx}
%   \end{acadtable}
\NewDocumentEnvironment{acadtable}{ O{htbp} m m }{%
  \begin{table}[#1]%
    \centering
    \captionof{table}{#3}\label{#2}%
}{%
  \end{table}%
}


% ─── Renombrar "Cuadro" → "Tabla" en español ─────────────────────────────────
\addto\captionsspanish{\renewcommand{\tablename}{Tabla}}


% ─── Nota al pie de tabla ────────────────────────────────────────────────────
\newcommand{\tablenote}[2]{%
  \par\vspace{2pt}%
  {\footnotesize\textsuperscript{#1}\,#2}%
}


%───────────────────────────────────────────────────────────────────────────────
% NOTA: tabularx NO puede envolverse en \NewDocumentEnvironment porque
% escanea el cuerpo literalmente buscando \end{tabularx}.
% Usar siempre directamente:
%   \begin{tabularx}{\textwidth}{l Y Z}   ← fuera de caja
%   \begin{tabularx}{\linewidth}{l Y Z}   ← dentro de tcolorbox
%───────────────────────────────────────────────────────────────────────────────
